\chapter{Features of NatQuery}
\label{ch:features}

This chapter lists what NatQuery offers, grouped by purpose. Chapter~\ref{ch:how} explains how each
mechanism works; Chapter~\ref{ch:bench} measures them.

\section{Feature Catalogue}

\begin{table}[htbp]
\centering
\caption{Feature catalogue.}
\label{tab:features}
\small
\renewcommand{\arraystretch}{1.15}
\begin{tabularx}{\textwidth}{L{2.4cm} L{3.7cm} Y}
\toprule
\hdr \textbf{Group} & \textbf{Feature} & \textbf{What it does} \\
\midrule
\multirow{5}{=}{Access control}
 & Pluggable authentication & Local accounts built in; LDAP, single sign-on or the ERP's own session via a provider interface \\
 & Policy-driven roles & Roles, inheritance and priorities in one YAML file; adding a role needs no code \\
 & Table and column grants & Hidden tables and columns are absent from the prompt and fail to resolve in the enforcer \\
 & Row filters & Per-user conditions bound from identity attributes, wrapped around every table reference \\
 & Masking & Columns can be shown masked (\code{redacted}, \code{last4}, \code{email}, \code{null}) but not filtered, sorted or joined on \\
\midrule
\multirow{5}{=}{SQL safety}
 & Never executes model SQL & The query is parsed, checked, rewritten and regenerated; only regenerated SQL runs \\
 & Statement allow-list & SELECT and set operations only; function allow-list; no cross joins, recursion or LATERAL \\
 & Bounded execution & Read-only transaction, timeout, row cap, join and nesting limits, always rolled back \\
 & Result guard & Pattern redaction (card numbers, tokens) and size caps on returned values \\
 & Input guard and limits & Injection screening (a speed bump, not the boundary), rate and model-call limits \\
\midrule
\multirow{4}{=}{Accuracy}
 & Repair loop & Fixable rejections go back to the model with a sanitised reason, at most twice \\
 & Schema selection & Weighted table ranking, foreign-key bridging and lookup tables for large schemas \\
 & Semantic hints & Table descriptions, synonyms and value hints for categorical columns \\
 & Wider-schema retry & An ``unanswerable'' answer on a trimmed schema is retried once with more tables \\
\midrule
\multirow{4}{=}{Databases}
 & PostgreSQL backend & Pooled connections, \code{SET LOCAL ROLE}, statement timeouts \\
 & Multi-schema mode & \code{schema.table} naming, \code{schema.*} grants, auto-completion of unambiguous names \\
 & Row-level security hook & Identity passed as transaction settings so database policies repeat the row scopes \\
 & Catalog-based extraction & Schema read from \code{pg_catalog}; the login needs no table privileges \\
\midrule
\multirow{4}{=}{Accountability}
 & Hash-chained audit log & Tamper-evident record of every request; optional HMAC key; \code{audit verify} command \\
 & Audit before execution & If the ``authorized'' record cannot be written, nothing runs \\
 & Opaque errors & Users see a generic message and a reference id; details only in the log \\
 & Secret redaction & Keys and passwords removed from logs; files are private to the owner (POSIX) \\
\midrule
\multirow{2}{=}{Operations}
 & Plan and cost analysis & \code{EXPLAIN}-based metrics and index suggestions for roles allowed to see plans \\
 & Command line and shell & Connect, query, administer schema, policy, users and audit from the terminal \\
\bottomrule
\end{tabularx}
\end{table}

\section{Extension Points}

Every external dependency sits behind an interface (Table~\ref{tab:seams}), so NatQuery can be adapted to
the real ERP without touching the security core.

\begin{table}[htbp]
\centering
\caption{Extension points.}
\label{tab:seams}
\small
\renewcommand{\arraystretch}{1.15}
\begin{tabularx}{\textwidth}{L{3.2cm} L{3.3cm} Y}
\toprule
\hdr \textbf{Concern} & \textbf{Built in} & \textbf{To replace it} \\
\midrule
Authentication & Local accounts (scrypt) & Implement \code{AuthProvider} and register it, or ship an entry point in the group \code{natquery.auth_providers} \\
Database & PostgreSQL & Implement \code{QueryBackend} and name its \code{sqlglot} dialect (for example Oracle) \\
Language model & Groq & Implement \code{LLMClient.complete}; any hosted or on-premises model \\
Input screening & Pattern and length guards & Implement \code{InputGuard} \\
Rate limiting, sessions & In-memory & Implement the limiter and session-store interfaces (for example Redis) \\
Roles and rules & YAML policy & Add a block; no code change \\
Masking strategies & Four built in & Add a template under \code{masking:} \\
\bottomrule
\end{tabularx}
\end{table}

\section{Command-Line Interface}

\begin{table}[htbp]
\centering
\caption{The \code{natquery} command.}
\label{tab:cli}
\small
\renewcommand{\arraystretch}{1.15}
\begin{tabularx}{\textwidth}{L{5.2cm} Y}
\toprule
\hdr \textbf{Command} & \textbf{Purpose} \\
\midrule
\code{natquery connect} & Configure the database and model for a workspace and extract the schema \\
\code{natquery select}, \code{reset}, \code{config} & Switch workspace, remove its configuration, show it with secrets masked \\
\code{natquery schema refresh} & Re-extract the schema (needs the permission \code{can_refresh_schema}) \\
\code{natquery schema hints} & Store example values of categorical columns for the prompt \\
\code{natquery policy init}, \code{check} & Create an example policy; validate the policy against the schema \\
\code{natquery user add}, \code{list}, \code{roles}, \code{attrs}, \code{passwd}, \code{disable}, \code{enable} & Manage local accounts, their roles and identity attributes \\
\code{natquery audit verify}, \code{tail} & Check the hash chain; show recent records \\
\code{natquery} & Interactive shell: log in, then ask questions in English \\
\bottomrule
\end{tabularx}
\end{table}

\section{Quality Assurance}

Because the enforcer is the security boundary, it has the heaviest testing (Table~\ref{tab:qa}).

\begin{table}[htbp]
\centering
\caption{How the code base is tested.}
\label{tab:qa}
\small
\renewcommand{\arraystretch}{1.15}
\begin{tabularx}{\textwidth}{L{4.2cm} Y}
\toprule
\hdr \textbf{Method} & \textbf{Detail} \\
\midrule
Unit and integration tests & 855 tests; branch coverage about 97\% with a 90\% gate; nothing excluded from coverage \\
Integration on real PostgreSQL & Enforcer corpus and full pipeline, write attempts that bypass the enforcer, statement timeouts, \code{SET LOCAL ROLE}, and 40 concurrent threads with isolation checks \\
Property-based fuzzing & Hypothesis generates queries against the enforcer; invariants are documented in the test file \\
Manual mutation testing & One security rule sabotaged at a time: 96 of 97 mutants killed; the survivor is an equivalent mutant \\
Pinned parser range & \code{sqlglot} is security-relevant; the declared range is the range the suite was run against \\
Three benchmark tracks & BIRD for accuracy, AdventureWorks for attacks, the ERP for scale and roles (Chapter~\ref{ch:bench}) \\
\bottomrule
\end{tabularx}
\end{table}

\FloatBarrier
